\section{Nonadic transports and memory preservation}
\label{sec:nonadic}

The reconstruction in Section~\ref{sec:memory} provides an
effective link between transport memory and the geometry of the
dodecaphasic record. The record generated by APP--TRIT--TPK determines its
incidence readings; the two compressions preserve, in separate
components, the terminal representation and its complements; combining
the memories with the charge restores the twelve coordinates. The geometric
direction and dimensional projector are obtained after this
restoration. The following development establishes the operator balance
underlying this chain and its behaviour under refinement,
spectral realization, and change of metric
\cite{HMT-VIII,hmtX}.

\subsection{Nine-phase calendar and orthogonal decomposition}

Let \(\mathcal H\) be a complex Hilbert space and let
\(C\in\mathcal B(\mathcal H)\). The orthogonal sum of nine copies,
\(\mathcal H^{\oplus9}\), separately preserves the nine positions
of the calendar. We define
\begin{equation}
 \begin{split}
 S_C(v_0,\ldots,v_8)&=(Cv_8,v_0,\ldots,v_7),\\
 J:\mathcal H\longrightarrow\mathcal H^{\oplus9},
 \qquad Jv&=\frac13(v,\ldots,v).
 \end{split}
 \label{nt:calendar}
\end{equation}
The normalization \(1/3=9^{-1/2}\) gives \(J^*J=I\). The operator
\(P=JJ^*\) is the orthogonal projection onto the uniform section;
its action replaces the nine components by their mean. Each
component crosses the junction containing \(C\) once during one
cycle, so that
\begin{equation}
 S_C^9=I_9\otimes C.
 \label{nt:full-return}
\end{equation}
The return of the calendar position thus realizes an action of
\(C\) on the memory fibre.

\begin{definition}
The uniform compression of the transport and its complement are
\begin{equation}
 T_C=J^*S_CJ=\frac{8I+C}{9},
 \qquad
 \eta_C=(I-P)S_CJ.
 \label{nt:compression}
\end{equation}
The first map has codomain \(\mathcal H\); the second
has codomain \(J(\mathcal H)^\perp\subset\mathcal H^{\oplus9}\).
\end{definition}

For \(z=(C-I)v\), the complement is written
\begin{equation}
 \eta_Cv=\frac1{27}(8z,-z,-z,-z,-z,-z,-z,-z,-z).
 \label{nt:complement}
\end{equation}
The sum of its components is zero. The coefficients of the
decomposition therefore arise from averaging eight
identical continuations and one transported continuation.

\begin{proposition}[Balance of compression and of the memory junction]
For every \(C\in\mathcal B(\mathcal H)\),
\begin{align}
 \eta_C^*\eta_C
   &=\frac8{81}(I-C)^*(I-C),\label{nt:eta-gram}\\
 I-T_C^*T_C
   &=\eta_C^*\eta_C+\frac19(I-C^*C).\label{nt:general-balance}
\end{align}
If \(C^*C=I\), the map
\(v\mapsto(T_Cv,\eta_Cv)\) is isometric.
\end{proposition}

\begin{proof}
Expression~\eqref{nt:complement} gives
\[
 \|\eta_Cv\|^2
 =\frac{8^2+8}{27^2}\|(C-I)v\|^2
 =\frac8{81}\|(C-I)v\|^2.
\]
Polarization proves~\eqref{nt:eta-gram}. On the other hand,
\[
 T_C^*T_C
 =\frac{64I+8C+8C^*+C^*C}{81}.
\]
Adding \eqref{nt:eta-gram} yields
\[
 T_C^*T_C+\eta_C^*\eta_C=\frac{8I+C^*C}{9}.
\]
Subtracting this result from \(I\) gives~\eqref{nt:general-balance}.
When \(C^*C=I\), the sum of the two Gram products is the
identity, which preserves all inner products.
\end{proof}

The term \(\eta_C^*\eta_C\) measures dispersion among the nine
components. The term \((I-C^*C)/9\) separately records
the change in norm at the memory junction. The unitary
realization places all visible variation in the first term and
preserves the norm in the full orthogonal sum.

\begin{proposition}[Preservation of an observable]
Let \(C\) be unitary and \(A=A^*\in\mathcal B(\mathcal H)\) satisfy
\(C^*AC=A\). Then
\begin{equation}
 A=T_C^*AT_C+
   \eta_C^*(I_9\otimes A)\eta_C
  =T_C^*AT_C+\frac8{81}(I-C)^*A(I-C).
 \label{nt:observable-balance}
\end{equation}
For \(A\succeq0\), both terms in the sum are positive.
\end{proposition}

\begin{proof}
The complement formula, applied to the sesquilinear product
defined by \(A\), proves the second equality. Expanding
the two terms gives
\[
 \frac1{81}\bigl(
 64A+8AC+8C^*A+C^*AC
 +8A-8AC-8C^*A+8C^*AC\bigr)
 =\frac{8A+C^*AC}{9}=A.
\]
Positivity is preserved under a map of the form
\(B\mapsto L^*BL\).
\end{proof}

\subsection{Iterating compression and preserving the fixed sector}

In this subsection we fix a unitary operator \(C\) and write
\(T=T_C\), \(\eta=\eta_C\). The transport \(S_C^N\) retains all
phases for \(N\) steps. The power \(T^N\), by contrast,
successively applies uniform compression; at each stage
\(\eta T^jv\) preserves the complement of the previous representation.
The following operators express this distinction.

\begin{proposition}[Isometric record at a finite horizon]
For every integer \(N\geq1\), the map
\begin{equation}
 V_Nv=(T^Nv,\eta v,\eta Tv,\ldots,\eta T^{N-1}v)
 \label{nt:finite-record}
\end{equation}
from \(\mathcal H\) to
\(\mathcal H\oplus(\mathcal H^{\oplus9})^{\oplus N}\)
is isometric. If \(A\) satisfies the hypotheses of
\eqref{nt:observable-balance}, then
\begin{equation}
 A=T^{*N}AT^N+
 \sum_{j=0}^{N-1}T^{*j}\eta^*(I_9\otimes A)\eta T^j.
 \label{nt:telescopy}
\end{equation}
\end{proposition}

\begin{proof}
Multiplying~\eqref{nt:observable-balance} by \(T^{*j}\) and
\(T^j\) gives
\[
 T^{*j}AT^j-T^{*(j+1)}AT^{j+1}
 =T^{*j}\eta^*(I_9\otimes A)\eta T^j.
\]
The telescoping sum proves~\eqref{nt:telescopy}. For \(A=I\),
the right-hand side is the Gram product of \(V_N\).
The step from \(N\) to \(N+1\) replaces only \(T^Nv\) by
\((T^{N+1}v,\eta T^Nv)\), preserving all memory
components previously produced.
\end{proof}

\begin{theorem}[Compression limit and unlimited memory]
Let \(P_{\mathrm{fix}}\) be the orthogonal projection onto
\(\ker(I-C)\). Then
\begin{align}
 \operatorname*{s\!-\!lim}_{N\to\infty}T^N
   &=P_{\mathrm{fix}},\label{nt:strong-limit}\\
 I&=P_{\mathrm{fix}}+
 \sum_{j=0}^{\infty}T^{*j}\eta^*\eta T^j.\label{nt:infinite-record}
\end{align}
The series converges strongly. For every bounded self-adjoint
observable satisfying \(C^*AC=A\), the following is also preserved:
\begin{equation}
 A=P_{\mathrm{fix}}AP_{\mathrm{fix}}+
 \sum_{j=0}^{\infty}
 T^{*j}\eta^*(I_9\otimes A)\eta T^j.
 \label{nt:infinite-observable}
\end{equation}
\end{theorem}

\begin{proof}
The spectral theorem for the unitary operator realizes \(C\) as
the coordinate function \(z\) on the unit circle.
Compression corresponds to \(t(z)=(8+z)/9\), and
\begin{equation}
 1-|t(z)|^2=\frac8{81}|1-z|^2,
 \qquad |z|=1.
 \label{nt:circular-defect}
\end{equation}
Thus \(t(z)^N\) converges to zero at every \(z\ne1\) and retains
the value one at \(z=1\). The bound \(|t(z)^N|\leq1\) allows
dominated convergence to be applied to each vector spectral measure.
This yields~\eqref{nt:strong-limit} and, by the same argument,
\(T^{*N}\to P_{\mathrm{fix}}\) strongly. The boundedness of
\(A\), together with \(\|T^N\|\leq1\), gives
\(T^{*N}AT^N\to P_{\mathrm{fix}}AP_{\mathrm{fix}}\).
The limit of~\eqref{nt:telescopy} proves both identities.
\end{proof}

The record
\[
 V_\infty v=(P_{\mathrm{fix}}v,\eta v,\eta Tv,\eta T^2v,\ldots)
\]
is therefore an isometry. The fixed sector preserves the
part that persists in the terminal representation; the other
components remain in the memory of the complements.

For the bilateral shift
\(C e_n=e_{n+1}\) on \(\ell^2(\mathbb Z)\), a fixed sequence
is constant. Square summability forces it to
vanish, so \(P_{\mathrm{fix}}=0\). The entire norm
is then represented in the complements. The spectrum of
this shift is the full circle, and the spectral
maximum of \(|t(z)|^N\) is one; hence
\(\|T^N\|=1\) for every \(N\). Strong convergence retains
this non-uniform behaviour.

For a cyclic permutation of \(d\) positions, the fixed sector
is the line of uniform vectors. If the permutation has
several orbits, each orbit contributes its own uniform direction.
This dimension determines the terminal part that must be preserved.

\begin{proposition}[Compression of an order-three transport]
Let \(C\) be unitary with \(C^3=I\). Then
\begin{equation}
 P_{\mathrm{fix}}=\frac{I+C+C^2}{3},
 \qquad
 T_C^*T_C=P_{\mathrm{fix}}+
 \frac{19}{27}(I-P_{\mathrm{fix}}).
 \label{nt:order-three}
\end{equation}
\end{proposition}

\begin{proof}
The average of the three powers is the orthogonal projector
onto the fixed vectors. Since \(C^*=C^2\),
\[
 T_C^*T_C=\frac{65I+8(C+C^2)}{81}
 =\frac{57I+24P_{\mathrm{fix}}}{81},
\]
which agrees with~\eqref{nt:order-three}.
\end{proof}

In the twelve-position realization on
\(\mathbb R^{12}\), \((Sx)_i=x_{i+1}\) with indices modulo
twelve, \(S^3\) has order four and \(S^4\) has order three.
The four orbits of \(S^4\) give
\(\operatorname{rank}P_{\mathrm{fix}}=4\).
The exceptional realization \(g_c\) of Article~X satisfies
\(g_c^3=I\) and \(\ker(I-g_c)=0\) on its own carrier;
there~\eqref{nt:order-three} reduces to
\(T_{g_c}^*T_{g_c}=19I/27\)
\cite{hmtX}. The coefficient and the fixed sector are transported
together with the domain.

The radial contraction \(M=qC\), \(0<q<1\), has another exact
record:
\begin{equation}
 L_Nv=\bigl(\sqrt{1-q^2}\,v,
 \sqrt{1-q^2}\,qCv,\ldots,
 \sqrt{1-q^2}\,q^{N-1}C^{N-1}v,q^NC^Nv\bigr).
 \label{nt:radial-record}
\end{equation}
Its Gram product is the identity because
\((1-q^2)\sum_{j=0}^{N-1}q^{2j}+q^{2N}=1\).
The terminal has norm \(q^N\|v\|\), and the infinite record
preserves the inner product. The factor \(q=5/9\) of the central
differential mode in Article~VIII applies to a full
cycle; its uniform distribution over nine steps uses
\(q^{1/9}S_C\), whose ninth power is
\(q(I_9\otimes C)\) \cite{HMT-VIII}.

\subsection{Record reduction and trace preservation}

The fixed direction of the complement in the nine components
allows a minimal realization in two summands. Define
\[
 D_C=\frac{\sqrt8}{9}(I-C).
\]
The preceding identities give
\(T_C^*T_C+D_C^*D_C=I\) for unitary \(C\).
The change of sign relative to the expression for
\(\eta_C\) corresponds to an orientation of the complementary
component and preserves its Gram product.

\begin{proposition}[Positive channel of compression with memory]
For unitary \(C\), the map
\begin{equation}
 \mathcal Q_C(X)=T_CXT_C^*+D_CXD_C^*
              =\frac{8X+CXC^*}{9}
 \label{nt:channel}
\end{equation}
is completely positive on \(\mathcal B(\mathcal H)\) and
preserves the identity. Its restriction to trace-class
operators preserves the trace.
\end{proposition}

\begin{proof}
Expanding the two terms in the first expression
cancels the products \(CX\) and \(XC^*\). The remaining
coefficients are \(72/81=8/9\) for \(X\) and \(9/81=1/9\)
for \(CXC^*\), proving~\eqref{nt:channel}.
For any integer \(m\geq1\), the amplification to matrices
of order \(m\) has the form
\[
 X\longmapsto
 (I_m\otimes T_C)X(I_m\otimes T_C)^*
 +(I_m\otimes D_C)X(I_m\otimes D_C)^*.
\]
Each summand preserves positivity; hence the map is
completely positive. The second expression in
\eqref{nt:channel} gives \(\mathcal Q_C(I)=I\).
If \(X\) is trace class, invariance of the trace under
unitary conjugation implies
\(\operatorname{Tr}\mathcal Q_C(X)=\operatorname{Tr}X\).
\end{proof}

The map \eqref{nt:channel} is also obtained by
preparing the isometry
\(v\mapsto T_Cv\otimes e_0+D_Cv\otimes e_1\) and taking
the partial trace over \(\mathbb C^2\).
Preservation of the full vector belongs to the isometry;
preservation of the trace of the reduced state belongs to
\(\mathcal Q_C\). Both statements follow from the same
orthogonal balance.

A complementary preservation appears in the nonadic
prolongation of finite algebras. For
\(\mathcal A_d=M_{9^d}(\mathbb C)\), with
\(\tau_d(a)=9^{-d}\operatorname{Tr}a\), the maps
\[
 \iota_d(a)=a\otimes I_9,\qquad
 \jmath_{d,j}(a)=a\otimes p_j,
 \quad p_j=|e_j\rangle\langle e_j|
\]
satisfy
\begin{align}
 \tau_{d+1}(\iota_d(a))&=\tau_d(a),
 &\iota_d(I)&=I,\label{nt:unital-trace}\\
 \tau_{d+1}(\jmath_{d,j}(a))&=\frac{\tau_d(a)}9,
 &\jmath_{d,j}(I)&=I\otimes p_j.\label{nt:corner-trace}
\end{align}
The proof uses
\(\operatorname{Tr}(a\otimes b)=
\operatorname{Tr}(a)\operatorname{Tr}(b)\),
\(\operatorname{Tr}I_9=9\) and \(\operatorname{Tr}p_j=1\).
The first inclusion preserves the entire fibre and the unit;
the second preserves a marked continuation and its support.

This distinction also appears in the tree of
enriched histories. Let \(\mathcal H_n\) be the surviving
prefixes of length \(n\), and let \(d(h)\geq1\) be the number of
retained outgoing emissions of \(h\). Each emission with multiplicity
is represented by distinct marked successors; the following sum runs over
this set of \(d(h)\) successors.
The weight law
\[
 \mu_{n+1}(h\epsilon)=\mu_n(h)\,p_h(\epsilon),
 \qquad p_h(\epsilon)=d(h)^{-1}
\]
preserves
\(\sum_{\epsilon}\mu_{n+1}(h\epsilon)=\mu_n(h)\).
Therefore, on
\(\mathscr L_n=L^2(\mathcal H_n,\mu_n)\), the pullback
\[
 R_nf(h\epsilon)=f(h)
\]
is isometric and preserves the unit function. Its adjoint is
\[
 E_ng(h)=\sum_{\epsilon}p_h(\epsilon)g(h\epsilon),
 \qquad E_nR_n=I,\quad E_n1=1.
\]
Indeed, grouping the sum for \(\|R_nf\|^2\) by predecessor
yields \(\sum_h\mu_n(h)|f(h)|^2\); the same grouping
in the inner product determines \(E_n\).
Memory and multiplicities remain present in the
domain of histories. These maps constitute the
functional realization of the refinement described in
Article~X \cite{hmtX}.

\subsection{Spectral chart and radius reciprocity}

The preceding transport admits spectral realizations
with their own domains. The M\"obius chart of VIII acts,
in finite coordinates, by
\begin{equation}
 \tau(z)=\frac{z-1}{z},
 \qquad z\in\mathbb C\setminus\{0\}.
 \label{nt:mobius}
\end{equation}
The inverse is \(z=(1-\tau)^{-1}\), for \(\tau\ne1\).
The extension to the Riemann sphere sends \(0\) to
\(\infty\) and \(\infty\) to \(1\).

\begin{proposition}[Spectral mirror and reciprocal radius]
Let \(z^\sharp=1-\overline z\). For
\(z\notin\{0,1\}\),
\begin{equation}
 \tau(z^\sharp)=\frac1{\overline{\tau(z)}}.
 \label{nt:mirror}
\end{equation}
If \(z=\sigma+it\), then
\begin{equation}
 |\tau(z)|^2
 =\frac{(\sigma-1)^2+t^2}{\sigma^2+t^2};
 \qquad
 |\tau(z)|=1\ \Longleftrightarrow\ \sigma=\frac12.
 \label{nt:critical-circle}
\end{equation}
\end{proposition}

\begin{proof}
Direct substitution gives
\[
 \tau(1-\overline z)
 =\frac{-\overline z}{1-\overline z}
 =\frac{\overline z}{\overline z-1}
 =\frac1{\overline{(z-1)/z}}.
\]
The ratio of squared moduli yields
\eqref{nt:critical-circle}; equality between numerator and
denominator is equivalent to \(-2\sigma+1=0\).
\end{proof}

Writing \(\tau(z)=r e^{i\theta}\), with \(r>0\), the
mirror preserves \(\theta\) and transforms \(r\) into \(r^{-1}\).
Its fixed set in this chart is the unit circle.
Conjugation \(z\mapsto\overline z\), in turn,
produces \(\tau\mapsto\overline\tau\) and reverses the angular
orientation. The two involutions preserve, respectively,
the angular coordinate and the radial coordinate.

On \(z=1/2+i\gamma\), with \(\gamma\in\mathbb R\),
\begin{equation}
 \tau_\gamma=\frac{\gamma+i/2}{\gamma-i/2},
 \qquad
 \gamma=\frac{i}{2}\frac{\tau_\gamma+1}{\tau_\gamma-1}
       =\frac12\cot\frac{\theta}{2},
 \quad \tau_\gamma=e^{i\theta}.
 \label{nt:spectral-phase}
\end{equation}
Compactification preserves the height through the inverse;
the point \(\tau=1\) corresponds to infinity.

Let \(\mathcal D=C_c^\infty(\mathbb R;\mathbb C)\), with its usual
test-function topology. The positive realization of the Weil
form explicitly uses its positive semidefiniteness, its
translation invariance, and its continuity in that topology, on the
domain proved by the corresponding spectral reader. If
\(\mathscr W(f,g)=\langle Ef,Eg\rangle\), its kernel is
\(\mathcal N=\{f\in\mathcal D:\mathscr W(f,f)=0\}\), and we define
\[
 \mathcal H_{\mathscr W}
 =\overline{\mathcal D/\mathcal N}.
\]
Translations induce a strongly continuous unitary
group \(U_t[f]=[f(\,\cdot-t)]\): invariance preserves the norm,
and continuity of translations in \(\mathcal D\) and of the form
proves strong continuity on a dense subspace and then on
the entire completion. With the convention
\(U_t=e^{itH}\), its self-adjoint generator acts on
test-function classes by \(H[f]=[if']\). The Cayley transform of
this realization is the bounded operator
\begin{equation}
 C_{\mathscr W}
 =(H+iI/2)(H-iI/2)^{-1}
 =I+i(H-iI/2)^{-1}.
 \label{nt:cayley-hilbert}
\end{equation}
Functional calculus makes it unitary, since the quotient
in~\eqref{nt:spectral-phase} has modulus one for
real \(\gamma\). The compression identity applies
to \(C_{\mathscr W}\) on \(\mathcal H_{\mathscr W}\).
The space, the domain of the generator, and the positivity
that permits its construction are the premises of this
spectral realization \cite{HMT-VIII}.

Bilateral memory, by contrast, acts on
\(\ell^2(\mathbb Z)\), and the permutations \(S^3,S^4\)
act on \(\mathbb R^{12}\). Each realization determines
its spectrum, its fixed sectors, and the action of one cycle.
The respective transport maps specify the
compositions between these spaces.

\subsection{Cayley on test functions and preservation of the arithmetic form}

The same arithmetic representation admits a Cayley reader
defined directly on test functions. Fix
\(a>b>1/2\) and
\[
 \mathscr E_b=
 \left\{f\in C^\infty(\mathbb R):
   \sup_x e^{b|x|}|f^{(j)}(x)|<\infty
   \text{ for every }j\geq0\right\}.
\]
Smooth compactly supported functions belong to
\(\mathscr E_b\). We define
\begin{align}
 (R_a^+f)(x)&=\int_0^\infty e^{-at}f(x+t)\,dt,
 &(R_a^-f)(x)&=\int_0^\infty e^{-at}f(x-t)\,dt,\label{nt:resolvents}\\
 C_af&=f-2aR_a^+f,
 &C_a^{-1}f&=f-2aR_a^-f.\label{nt:test-cayley}
\end{align}
Each derivative of \(R_a^\pm f\) satisfies the bound for the
corresponding seminorm of \(f\), multiplied by
\((a-b)^{-1}\). This follows from
\(e^{b|x|}|f(x\pm t)|\leq M_0e^{bt}\).
Both maps preserve \(\mathscr E_b\).

With the transform
\(\widehat f(\xi)=\int_{\mathbb R}f(x)e^{-i\xi x}\,dx\),
Fubini's theorem gives
\begin{equation}
 \widehat{C_af}(\xi)
 =c_a(\xi)\widehat f(\xi),\qquad
 c_a(\xi)=\frac{\xi-ia}{\xi+ia}.
 \label{nt:cayley-multiplier}
\end{equation}
The multiplier of the second formula
in~\eqref{nt:test-cayley} is \(c_a^{-1}\), proving
the inverse. For \(\xi\in\mathbb R\),
\(|c_a(\xi)|=1\); Plancherel provides its unitary
extension to \(L^2(\mathbb R)\).

To specify arithmetic preservation, let
\[
 \mathcal C_{f,g}(u)
 =\int_{\mathbb R}f(x)\overline{g(x-u)}\,dx,
 \qquad
 P_f^\pm=\int_{\mathbb R}f(x)e^{\pm x/2}\,dx.
\]
The lengths \(\ell_{p,k}=k\log p\) and weights
\(w_{p,k}=(\log p)p^{-k/2}\) are the prime--power
representations of the arithmetic reader. The complete form
of Article~VIII is written
\begin{equation}
 \begin{split}
 \mathscr W(f,g)={}&c_\Gamma\mathcal C_{f,g}(0)\\
 &+\int_0^\infty\frac{e^{-u/2}}{1-e^{-2u}}
 \bigl[2\mathcal C_{f,g}(0)
       -\mathcal C_{f,g}(u)-\mathcal C_{f,g}(-u)\bigr]\,du\\
 &-\sum_{p,k}w_{p,k}
   \bigl[\mathcal C_{f,g}(\ell_{p,k})
        +\mathcal C_{f,g}(-\ell_{p,k})\bigr]
 +P_f^+\overline{P_g^-}+P_f^-\overline{P_g^+},
 \end{split}
 \label{nt:weil-complete}
\end{equation}
where \(c_\Gamma=\psi(1/4)-\log\pi\), and
\(\psi=\Gamma'/\Gamma\) is the logarithmic derivative of the
gamma function. These functions and constants enter
the subsequent analytic realization of the arithmetic
reader of the corpus.

The bound
\[
 |\mathcal C_{f,g}(u)|
 \leq M_fM_g\bigl(|u|+b^{-1}\bigr)e^{-b|u|}
\]
controls the prime sum by a convergent series of the type
\(\sum_{n\geq2}(\log n)(1+\log n)n^{-b-1/2}\).
The symmetric difference of correlations is \(O(u^2)\)
at the origin, while the gamma weight is \(O(u^{-1})\).
At infinity that weight decays exponentially.
The margin \(b>1/2\) also ensures convergence of
the two polar readers. This fixes the domain
of all terms in~\eqref{nt:weil-complete}.

\begin{proposition}[Arithmetic invariance and nonadic balance]
For \(f,g\in\mathscr E_b\),
\begin{equation}
 \mathscr W(C_af,C_ag)=\mathscr W(f,g).
 \label{nt:weil-invariance}
\end{equation}
With \(\mathcal T_a=(8I+C_a)/9\), we have
\begin{equation}
 \mathscr W(f,g)
 =\mathscr W(\mathcal T_af,\mathcal T_ag)
 +\frac8{81}\mathscr W((I-C_a)f,(I-C_a)g).
 \label{nt:weil-balance}
\end{equation}
\end{proposition}

\begin{proof}
The multiplier \(c_a\) is unitary on the real line and
commutes with translations. Consequently,
\(\mathcal C_{C_af,C_ag}(u)=\mathcal C_{f,g}(u)\)
for every \(u\). By Fubini, the polar readers transform
according to
\[
 P_{C_af}^+
 =-\frac{a-1/2}{a+1/2}P_f^+,\qquad
 P_{C_af}^-
 =-\frac{a+1/2}{a-1/2}P_f^-.
\]
Their factors are real and reciprocal, so that
each cross polar product remains unchanged.
Combining the terms of
\eqref{nt:weil-complete} proves~\eqref{nt:weil-invariance}.
The sesquilinear expansion of the right-hand side of
\eqref{nt:weil-balance} cancels the mixed terms and
yields
\(\bigl(8\mathscr W(f,g)+
\mathscr W(C_af,C_ag)\bigr)/9\), equal to
\(\mathscr W(f,g)\).
\end{proof}

Finite iteration preserves the complete form:
\[
 \mathscr W(f,g)
 =\mathscr W(\mathcal T_a^Nf,\mathcal T_a^Ng)
 +\frac8{81}\sum_{j=0}^{N-1}
 \mathscr W((I-C_a)\mathcal T_a^jf,
            (I-C_a)\mathcal T_a^jg).
\]
This sesquilinear identity simultaneously preserves
gamma, primes, and polar terms. Its evaluation in a
positive realization allows an interpretation as an energy
distribution; its algebraic validity precedes that
interpretation. The margin \(a>b>1/2\) belongs to the
exponential test-function domain of this construction.
The Cayley transform~\eqref{nt:cayley-hilbert} with parameter \(1/2\)
belongs to the Hilbert-space domain of its generator.

\subsection{Metric covariance of the dodecaphasic flow}

The reconstruction in Section~\ref{sec:memory} determines
\(u_K\in\mathbb R^{12}\), with
\(\mathbf1^{\mathsf T}u_K=0\) and \(u_K\ne0\).
We write \(A_K=\operatorname{diag}(u_K)\), with the direction and
normalization declared in \eqref{eq:u-main}, and
\[
 g_s=e^{sA_K}\in\operatorname{GL}(12,\mathbb R).
\]
The zero trace of \(A_K\) implies \(\det g_s=1\).
Each \(g_s\) transports the weights and geometric
coordinates according to the corresponding realization.
Preserving the memory balance between charts
also requires transporting the inner product.

\begin{proposition}[Metric transport of compression and complement]
Let \(V\) be a finite-dimensional Hermitian space, \(C_0\) unitary,
and \(g:V\to V\) invertible. We write
\(g^{-*}=(g^{-1})^*\) and define
\begin{equation}
 C_g=gC_0g^{-1},\qquad
 G_g=g^{-*}g^{-1},\qquad
 \langle v,w\rangle_g=\langle G_gv,w\rangle.
 \label{nt:metric-chart}
\end{equation}
Then \(C_g^*G_gC_g=G_g\), and
\begin{align}
 T_{C_g}g&=gT_{C_0},\label{nt:metric-naturality}\\
 \eta_{C_g}g&=g^{\oplus9}\eta_{C_0},\label{nt:eta-naturality}\\
 G_g&=T_{C_g}^*G_gT_{C_g}
 +\frac8{81}(I-C_g)^*G_g(I-C_g).\label{nt:metric-balance}
\end{align}
\end{proposition}

\begin{proof}
The positivity of \(G_g\) follows from
\(\langle G_gv,v\rangle=\|g^{-1}v\|^2\).
Substituting~\eqref{nt:metric-chart},
\[
 C_g^*G_gC_g
 =g^{-*}C_0^*g^*g^{-*}g^{-1}gC_0g^{-1}
 =g^{-*}C_0^*C_0g^{-1}=G_g.
\]
The identity \(C_gg=gC_0\) gives
\eqref{nt:metric-naturality} upon adding \(8g\) and dividing
by nine. It also gives
\((C_g-I)g=g(C_0-I)\); substituting it into
\eqref{nt:complement} proves~\eqref{nt:eta-naturality}.
Finally, the expansion of
\eqref{nt:observable-balance}, with \(A=G_g\) and
\(C_g^*G_gC_g=G_g\), proves
\eqref{nt:metric-balance}.
\end{proof}

Applying the proposition to \(g=g_s\) and to the operators
\(C_0=S^3,S^4\) of the dodecaphasic realization transports
the compressions, the memories, and their inner products
through explicit commutative squares.
The identification \(g_s:(V,\langle\cdot,\cdot\rangle)
\to(V,\langle\cdot,\cdot\rangle_{g_s})\) is isometric.
The formula therefore also transports the fixed
sectors and the records at every depth:
\[
 P_{\mathrm{fix},s}
 =g_sP_{\mathrm{fix},0}g_s^{-1}.
\]
The projector on the right-hand side is orthogonal with respect
to \(G_{g_s}\).

The volume condition and the metric condition have
distinct effects. For the rational matrix
\[
 g=\begin{pmatrix}2&0\\0&1/2\end{pmatrix},
 \qquad \det g=1,
\]
direct substitution into the compression gives
\[
 \frac{8I+g}{9}
 =\begin{pmatrix}10/9&0\\0&17/18\end{pmatrix},
 \qquad
 T_g^*T_g+\frac8{81}(I-g)^*(I-g)
 =\frac{8I+g^*g}{9}.
\]
The first coordinate expands in the Euclidean
metric. The balance~\eqref{nt:metric-balance},
by contrast, preserves the metric unit when a
unitary operator \(C_0\) is transported together with
its inner product. The determinant controls
volume; \(G_g\) controls norms and cross
products.

The flow satisfies \(g_{-s}=g_s^{-1}\).
The radius reciprocity in
\eqref{nt:mirror} belongs to the spectral chart,
and inversion of \(g_s\) belongs to the
dimensional realization of the record. The constructive
link used here preserves the
explicit maps from memory to \(K\) and from
\(K\) to its geometric direction, followed by
the metric transport of each realization.
This composition maintains, within a single development,
exact reconstruction, positive refinement,
boundary recording, and preservation
of inner products.
